\section{Problems and results}
\label{sec:results}

The paper contains a variety of different but related results involving classical tiling and
quantum Hamiltonian problems with translational invariance.  In this section, we will
summarize the different variants, giving the proofs and more detailed discussion of
each variant in later sections.  While there are certain recurring techniques, the
details of the different cases vary considerably.

\subsection{Tiling results}

\begin{definition}{\bf \tiling}

\noindent
{\bf Problem Parameters:} A set of tiles $T = \{t_1,\ldots,t_m\}$.
A set of horizontal constraints $H \subseteq T \times T$ such that if
$t_i$ is placed to the left of $t_j$, then it must be the case that
$(t_i,t_j) \in H$. A set of vertical constraints $V \subseteq T \times T$ such that if
$t_i$ is placed below $t_j$, then it must be the case that
$(t_i,t_j) \in V$. A designated tile $t_1$ that must be placed in the four corners of the grid.


\noindent
{\bf Problem Input:} Positive integer $N$, specified in binary.



\noindent
{\bf Output:} Determine whether there is a valid tiling of an $N \times N$ grid.
\end{definition}





\begin{theorem}
\label{th:unweighted-tiling}
There exists a $(T,H,V,t_1)$ for which \tiling\ is \NEXP-complete.
\end{theorem}


A similar comment holds for all the hardness results given in this paper.
Although \expref{Theorem}{th:unweighted-tiling} is subsumed by a result proven
independently in~\cite{BGL09}, for completeness,
we give the proof in \expref{Section}{sec:tiling}. The basic idea is that the
corner tiles are used to create a border around the perimeter of the grid which
allows us to implement special rules at the top and bottom rows.
The interior of the grid is tiled in two layers.
The set of colors for the interior tiles is the cartesian product of the
layer 1 colors and the layer 2 colors.
The rules for the two layers are independent except at the boundary.
Thus, in the interior of the grid, two tiles can be placed next to each
other as long as the placement is consistent with  the layer 1 colors for
the tiles (according to the layer 1 rules) and the layer 2 colors for the tiles
(according to the later 2 rules).
In the construction, each layer implements
the action of a Turing machine.
The first TM proceeds from top to bottom
on layer 1 and the second proceeds from bottom to top on layer 2.
The first TM takes no input and acts
as a binary counter for $N$ steps. The bottom row
of  the first layer then
holds a binary number that is $\Theta(N^{1/k})$ for some constant $k$.
The rules for the lower boundary are then used
to copy the output from the binary counter to the bottom row of layer 2, which
acts as the input to a non-deterministic Turing machine which computes a \NEXP-complete
language.
The rules for the top boundary check whether
the final configuration on layer 2 is an accepting  state.
\cite{BGL09} goes about their proof somewhat differently although the end result is the
same. They look at a generalization of strings called \emph{pictures} which are essentially matrices
over a finite alphabet. \emph{Picture languages} are sets of these objects.
In particular, they study the complexity of recognizing unary picture languages
whose elements are squares.
They show that such a language corresponds to the set of tilable squares according to a finite
set of tiling rules if and only if the language defined by the binary representation
of the sizes of the squares is in $\NEXP$ time. This proof uses a binary counter similar to our
construction. In a separate proof, they show that finite tiling systems can
simulate an arbitrary Turing Machine.
It turns out that the width of the grid is not important in this
construction as long as it is an integer
at least $N$.

Note that it is important that we chose to have the input $N$ provided in binary.
If it were instead given in unary, there would only be one instance per problem size,
and the problem would be trivially in $\Pclass$/poly.  Thus, in order to prove
a meaningful hardness result, we are forced to move up the exponential hierarchy
and prove the problem is $\NEXP$-complete rather than $\NP$-complete.
